% =============================================================================
% Chapter 8 -- BCS Criteria and Self-Reflection
% =============================================================================
\chapter{BCS Criteria and Self-Reflection}
\label{ch:bcs}

\section{BCS Learning Outcomes Demonstrated}

The British Computer Society accreditation framework requires the
dissertation to demonstrate the following six learning outcomes.
Each item below is addressed in the named chapter or section.

\begin{enumerate}[label=\textbf{B\arabic*.},leftmargin=2.4em]
    \item \textbf{Application of practical and analytical
        skills.}\quad Practical skills are demonstrated by the
        end-to-end implementation in \Cref{ch:impl}, which
        adapts a published GPU-targeted research codebase to a
        consumer CPU host without changing the comparison
        semantics. Analytical skills are demonstrated by the
        diagnostic chain in \Cref{sec:fishernoise}, which moves
        from a quantitative gap (\Cref{tab:avg_perf}) through a
        within-task drift control (\Cref{tab:forgetting}) to a
        mechanistic hypothesis (Fisher on noise) that is itself
        framed as a falsifiable prediction for future work.

    \item \textbf{Innovation and creativity.}\quad The
        innovation is methodological rather than algorithmic:
        the dissertation re-frames a published continual
        learning benchmark as a study of the implicit compute
        assumptions in that benchmark, and reports an
        \emph{inversion} of the published ranking under those
        assumptions. The Fisher-on-noise diagnosis is, to the
        author's knowledge, not stated explicitly in the
        existing CL or COOM literature.

    \item \textbf{Synthesis of information, ideas and
        practices.}\quad \Cref{sec:background} integrates three
        threads --- regularisation-based continual learning,
        the COOM benchmark, and the deep RL plasticity
        literature --- into a single argument for studying the
        low-budget regime. The plasticity-loss work
        \citep{lyle2022plasticity, nikishin2022primacy} is the
        external thread that motivates the Fisher-on-noise
        hypothesis; without it the diagnosis would be
        unmotivated.

    \item \textbf{Real need in a wider context.}\quad On-device,
        online, and student-scale continual reinforcement
        learning deployments cannot afford the per-task compute
        budgets used by published benchmarks. This dissertation
        surfaces a previously under-discussed failure mode in
        that operational regime, and the negative-result
        framing in \Cref{sec:wider} argues directly that the
        finding is relevant to practitioners.

    \item \textbf{Self-management of the project.}\quad Project
        management is reflected in the overnight queue
        infrastructure (\Cref{sec:queue}), the up-front design
        decisions documented in \Cref{ch:design}, and the
        handling of the most significant unplanned change ---
        the decision to commit to a negative-result framing
        once the data was in --- which is discussed in the
        self-reflection below.

    \item \textbf{Critical self-evaluation.}\quad The
        first-person section that follows addresses this
        criterion explicitly.
\end{enumerate}

\section{Self-Reflection}
\label{sec:reflection}

I began this project expecting to reproduce a published
continual reinforcement learning baseline on commodity
hardware, and I ended it with a finding that the published
baselines do not transfer to that hardware regime in the simple
way I had assumed. The shift from \emph{replication} to
\emph{negative-result study} was the single most important
course-correction in the project, and it took me longer than it
should have to commit to it.

The first weeks of the project were spent on infrastructure: a
working CPU TensorFlow build of the COOM training loop, an
overnight experiment queue, a TSV log loader. I expected, when
the first end-to-end run finished, to see a CO4 success curve
that looked qualitatively like the COOM paper's headline
figures. What I saw instead was a Fine-Tuning curve that
looked broadly reasonable and an EWC curve that flatlined on
the final task. My instinct was to assume an implementation
bug. I spent something like a week chasing one --- inspecting
the Fisher computation, varying the regularisation coefficient,
checking the replay buffer reset behaviour, comparing the
training loop step-by-step against the upstream code. Nothing
came back as broken. The regularisation-loss curve in
\Cref{fig:lossreg} was, by then, the strongest piece of
evidence I had that the implementation was correct: the EWC
penalty fires at exactly the right step counts, with exactly
the right shape.

The hardest moment of the project was the conversation with
myself in which I had to accept that the difference from the
paper was not a bug but a finding. Two things made the
acceptance slow. The first was sunk cost: I had told myself
the project was a replication, and turning a replication into
a negative-result study felt, at the time, like admitting I
had failed at the original goal. The second was a more
interesting form of self-doubt: I worried that an undergraduate
finding that the published benchmark did not hold under
constrained compute would read as either naive or arrogant,
depending on the reader. What changed my mind was rereading
the COOM paper carefully and noticing that it does not, in
fact, claim its results generalise to lower compute --- the
absence of such a claim is itself an opening for the question
this dissertation now asks. The negative-result framing is
the honest one.

If I were to start the project again, three things would be
different. First, I would invest in the
overnight-queue infrastructure on day one rather than during
the second sprint; the queue is what made the seed band on
EWC affordable, and without it the seed-variance argument in
\Cref{sec:results} would be much weaker. Second, I would
budget multiple seeds for every method from the start rather
than backfilling them only for EWC; the partial satisfaction
of \ref{req:variance} is the largest concrete weakness of the
study and is entirely a planning failure. Third, I would
commit earlier to a small but well-instrumented experiment
(CO4, four methods) over a large but thinly-instrumented one
(CO8, all methods, no seed band). The former is what I ended
up reporting; the latter is what I drifted into in the first
half of the project.

What I carry forward is a sharper instinct for when to stop
debugging and start trusting the data. The Fisher-on-noise
diagnosis was visible in the regularisation-loss plot for
weeks before I named it. Empirical research, I now think, is
mostly about giving yourself permission to draw the
conclusions your measurements are already pointing at, and the
reflexive ``it must be a bug'' response is the single most
expensive habit to unlearn. My view of the continual-learning
literature has also matured: published rankings are
conditional on the experimental setup that produced them, and
in a field where setups vary by orders of magnitude in
compute, the ranking is best treated as one data point about
that setup rather than as a property of the methods
themselves. That is the lesson I want my own future work to
carry.
